\section{Missing studies, missing results and selective outcome reporting}
\label{sec:p15-missing-evidence}

A synthesis can be precise for the evidence it contains while omitting
evidence that would change its conclusion. The relevant denominator is
the eligible research, not just the reports found or the results that
can be entered into a calculation. This treatment uses Chapter~13 of
the Cochrane Handbook, version~6.5, last updated August~2024,
Sections~13.1--13.4 \cite{PageHigginsSterne2024Cochrane13}.
The pooling models and within-study missing-participant problem remain
those of \cref{sec:p15-evidence-synthesis,def:p15-synthesis-missing}.

\begin{hbdefinition}{def:p15-missing-evidence-units}{Study, outcome, result and availability}
A study is the research undertaking; several reports can describe it.
A participant outcome is defined by its domain, measurement instrument,
metric and time point. A study result additionally specifies its
aggregation, intervention/comparator contrast and analysis population
or model, and supplies an effect estimate with a precision measure,
or the statistics needed to calculate both.

For a specified synthesis, distinguish:
\begin{enumerate}
\item An unidentified eligible study, for which even its existence
      may be unknown to the reviewers.
\item An identified study with no usable eligible result, although
      a register entry, protocol or report establishes the study.
\item A partially available result: for example a point estimate
      without its variance, an arm mean without the comparator,
      or only a statement about statistical significance.
\item A reported result selected from alternative outcomes, times,
      metrics or analyses because of its findings.
\end{enumerate}
The fourth can bias an available study result without making a whole
study missing from a broadly defined meta-analysis. Missing participant
outcomes within a study are another problem. A study that did not
measure the specified outcome is not a study that generated its result
and withheld it. An unknown measurement status does not establish either
case \cite[secs.~13.1, 13.3.1.2--13.3.1.3]{PageHigginsSterne2024Cochrane13}.
\end{hbdefinition}

\begin{hbexample}{def:p15-missing-outcome-eligibility}{Availability depends on the declared result}
This is a stipulated example. A trial measured the same symptom scale
at two and six weeks. It reports the two-week difference and precision,
but withholds the six-week comparison because of its findings.
For a synthesis specifically targeting six weeks, its eligible result
is missing. For a predeclared synthesis permitting any measurement
through eight weeks, the two-week result is available and eligible.
Choosing it because the six-week result was less favourable can still
give selection bias in the reported result. The broader outcome rule
has not cured that bias, and broadening the rule after seeing the
findings would add a review-level selection decision.
\end{hbexample}

\begin{hbdefinition}{def:p15-missing-selection}{A finite result-dependent availability model}
Let \(Y\) be a numerical result on the declared scale, with finite
support \(\mathcal Y\) and probabilities \(p(y)\).
Let \(R=1\) mean that it becomes available to the synthesis, and
\(s(y)=\Pr(R=1\mid Y=y)\). Require
\(0\leq s(y)\leq1\) and
\(\rho=\sum_y p(y)s(y)>0\). This is a stipulated selection
model, not an estimated explanation of a particular review.
The conditional available-result distribution is \(q(y)\).
Both a report-level selection rule and an outcome-level reporting
rule can produce such a distribution. The numerical result can
be enlarged to a vector including standard error, time or other
selection-relevant attributes when those matter.
\end{hbdefinition}

\begin{hbtheorem}{thm:p15-missing-selection}{proved=annex:p15.missing.selection}{How availability changes a distribution and what it cannot identify}
Under the preceding model,
\[
 q(y)=\frac{p(y)s(y)}\rho,\qquad
 E(Y\mid R=1)-E(Y)=\frac{\operatorname{Cov}(Y,s(Y))}\rho.
\]
If all \(s(y)>0\) are specified, the full law is reconstructed as
\[
 p(y)=\frac{q(y)/s(y)}{\sum_xq(x)/s(x)},\qquad
                  E\!\left[\frac{RY}{s(Y)}\right]=E(Y).
\]
Neither result identifies \(s\) from the observed law alone.
With equally likely full results \(-1,1\), and availability
\(s(-1)=1/3,s(1)=1\), the available mean is \(1/2\) although
the full mean is zero. Exactly the same available distribution
arises with no selection and full probabilities \(1/4,3/4\),
whose full mean is \(1/2\). Results whose selection probability
is zero cannot be recovered by inverse weighting.
\end{hbtheorem}

Selection can depend on magnitude, direction, a significance threshold,
time to dissemination, publication language or venue, and subsequent
indexing or citation. The investigator, sponsor, journal and review
search can contribute different stages. Recording these stages helps
separate a missing report from a report with a missing outcome.
Direction need not always favour the experimental intervention; an
investigator with the opposite expectation can suppress favourable
findings. A missing result does not establish the motive for its absence.

\section{Search, registries, protocols and inception cohorts}
\label{sec:p15-missing-search}

\begin{hbmethod}{def:p15-missing-search}{Recover studies and results with a retained accounting}
Search appropriate general and specialist bibliographic sources,
prospective study registers and results registers, conference and
other non-journal sources, and relevant regulatory or sponsor records.
Look for study protocols, statistical analysis plans and their dated
amendments. Link all reports to one study identity. A results register
can supply numerical findings from a study that has no journal article;
a protocol or a registration alone normally supplies plans rather than
the missing effect and precision.

For regulated interventions, seek clinical study reports and regulatory
assessment documents where accessible. They can identify otherwise
unseen trials or outcomes, but access, redaction, large document sets,
incompatible analyses and conflicting versions need to be accounted for.
Request missing results from investigators or sponsors; preserve the
request, response and extraction provenance. Treat data supplied by an
interested party as evidence requiring verification, not a declaration
that its entire study portfolio has been supplied.

Compare what was planned, what was actually measured, what was analysed,
and what is reported. Use original registry versions as well as later
updates. A legitimate unannounced protocol amendment is possible;
pre-specification alone does not prove the measurement occurred.
If plans are unavailable, compare methods with results, earlier abstracts
with later papers, and outcomes that are necessarily or commonly
measured together. Ask for an explanation before equating absence with
selective withholding. Keep every otherwise eligible study in the
availability record even if it cannot enter the numerical calculation
\cite[secs.~13.2.1, 13.3.1.3]{PageHigginsSterne2024Cochrane13}.
\end{hbmethod}

An inception cohort identifies its studies before their results are
known. A prospectively registered cohort, a complete grant or regulatory
frame, a fully enumerated sponsor portfolio, or a prospectively planned
set of trials can therefore make the study denominator known.
This requires completeness of that declared frame and its enumeration.
It does not make all outcomes or results available, and it does not
identify unregistered research outside the frame. Restricting a review
to that cohort changes its eligibility boundary, with possible losses
of precision and applicability. Registration adoption differs across
periods and research settings; a current register is not a census of
all historical research \cite[sec.~13.2.2]{PageHigginsSterne2024Cochrane13}.

\begin{hbtheorem}{thm:p15-missing-cohort}{proved=annex:p15.missing.cohort}{A complete cohort and the target it actually covers}
Let a finite full target assign fixed positive weights \(b_i\) to
studies, with \(\sum_i b_i=1\), and let a nonempty inception
cohort \(G\) have total weight \(\lambda=\sum_{i\in G}b_i\).
If its enumerated study list is complete, no study eligible under that
cohort-only rule is unidentified. Its results may still be missing.
Writing \(\theta_G\) for its normalized weighted effect, when
\(0<\lambda<1\) the full target is
\[
                \theta=\lambda\theta_G+(1-\lambda)\theta_{G^c}.
\]
If only \(L\leq\theta_i\leq U\) is known outside the cohort,
the full target lies in
\[
 [\lambda\theta_G+(1-\lambda)L,
                  \lambda\theta_G+(1-\lambda)U],
\]
with both extremes attainable. At \(\lambda=1\) the targets
coincide. A complete cohort alone does not identify the full target
when the excluded studies have unspecified effects and positive weight.
These are target-accounting statements, not confidence intervals.
\end{hbtheorem}

\paragraph{A documented search example.}
Eyding and colleagues' 2010 review sought double-blind randomized trials
of at least six weeks of acute treatment in adults with major depression,
comparing reboxetine with placebo or SSRIs. Searches through February~2009
included bibliographic sources, trial registers and regulatory websites;
the authors also sought previously unpublished manufacturer data.
They analysed 13 eligible acute-treatment trials with 4098 patients,
of whom 3033, about 74 percent, belonged to unpublished data. This is
a patient-data denominator, not a percentage of trials.
Some papers reported only subgroups or omitted the main outcomes, so
the existence of an article did not establish availability of the
eligible full-trial results. Their comparison of published with combined
published/unpublished evidence found overestimated benefits and
underestimated harms in the published subset. This is a documented
comparison for that review and its outcomes, not a claim that every
unpublished study has an unfavourable result or that searching a registry
always changes the estimate \cite[Methods, Study selection and Publication bias]{EydingEtAl2010MissingEvidence}.

\section{Risk-of-bias assessment for missing evidence}
\label{sec:p15-missing-robme}

The selected chapter's ROB-ME assessment concerns one specified
pairwise meta-analysis result. It can be used even when only one
identified study has a result available for that calculation.
It is not the separate ROB-MEN network tool, a general qualitative
synthesis instrument, or the study-level assessment of selection of
the reported result. When no summary meta-analytic result is produced,
the complete selected ROB-ME procedure does not apply, but missing
evidence still belongs in the interpretation
\cite[sec.~13.3]{PageHigginsSterne2024Cochrane13}.
The precise operating presentation used below is the tool's
1 October 2023 crib sheet and template \cite{ROBME2023Cribsheet}.

\begin{hbdefinition}{def:p15-missing-results-matrix}{The study-by-synthesis availability record}
For every identified study meeting the other eligibility criteria,
record its identifiers, all consulted sources, participants analysed
(or allocated if the analysed number is unknown), and one classification
for each specified meta-analysis:
\[
 \begin{aligned}
 \checkmark\ &:\text{usable eligible result available};\\
 \sim\ &:\text{unavailable for a reason unrelated to its findings};\\
 ?\ &:\text{eligible result generation unclear};\\
 X\ &:\text{unavailable likely because of its findings}.
 \end{aligned}
\]
For each entry retain its explanation and source. Include otherwise
eligible studies without published reports and those excluded from
the calculation solely because data were unusable. Distinguish an
outcome not measured from one whose measurement status is uncertain.
Record any known sign, significance information or partial summary
without replacing the missing numerical result by zero.
\end{hbdefinition}

\begin{hbmethod}{def:p15-missing-robme-process}{The selected ROB-ME planning and questions}
\begin{enumerate}
\item Select the important meta-analyses for assessment, ideally in
      advance. Define their populations, intervention/comparator,
      eligible study designs, outcome domain, instrument, metric,
      aggregation, time points and analysis methods. A change in that
      definition can change whether a result is missing.
\item Assemble the availability matrix from every identified eligible
      study and investigate explanations. Results described only as
      nonsignificant, pooled across arms, or incompletely reported
      remain missing for a calculation needing effect and precision.
      A measurement failure, an outcome genuinely not collected, or
      an accidentally absent appendix can be unrelated to the findings;
      document the evidence for that explanation.
\item Assess potential unidentified studies for the review as a whole.
      Planning question 3.1 asks whether eligibility is restricted to
      a genuine inception cohort. If not, question 3.2 asks whether
      all eligible studies should be identifiable independently of
      their results; if yes or probably yes, 3.3 asks whether the search
      likely found them all. A negative or probably negative answer
      to either relevant latter question marks potential missing studies.
      Here an affirmative inception-cohort answer is protective; it is
      not a count of high-risk responses.
\item For each meta-analysis, record its estimate and interval, model,
      included studies and participants. Answer the following eight
      questions, preserving the tool's applicability conditions.
\end{enumerate}

Questions 4.1--4.4 address identified studies:
\begin{enumerate}
\item[4.1] Is any matrix entry \(X\), indicating an unavailable
             result likely related to its findings? Answer yes or no.
\item[4.2] If yes, would adding those results probably change the
             summary magnitude notably? Consider their possible weight,
             number, heterogeneity and any known direction.
\item[4.3] Is any matrix entry \(?\), leaving generation of an
             eligible result unclear? Answer yes or no.
\item[4.4] If yes, could that potentially missing amount probably
             change the summary notably?
\end{enumerate}
Questions 4.5--4.8 add unidentified studies and the observed pattern:
\begin{enumerate}
\item[4.5] Does the review-wide planning indicate possible
             result-related failure to identify studies? Answer yes or no.
\item[4.6] If yes, were those studies likely to have generated results
             eligible under this meta-analysis's outcome definition?
\item[4.7] If any of 4.1, 4.3 or 4.5 is yes, does the observed pattern
             suggest systematically different missing results, after
             considering competing explanations for small-study effects?
\item[4.8] If any of 4.2, 4.4, 4.6 or 4.7 is yes, probably yes or
             no information, do plausible sensitivity analyses indicate
             notable change due to missing results?
\end{enumerate}
For 4.2 and 4.4 retain yes, probably yes, probably no, no,
no information and not applicable. The tool reserves no information
there for an unclear size of the missing or potentially missing study.
For the other conditional questions use yes, probably yes, probably
no, no and not applicable. Use not applicable to skip an untriggered
question; 4.7 can also be not applicable when too few results exist
to assess a pattern. Probably responses have the same grading effect
as their corresponding firm response, but retain the weaker evidence.
The selected tool permits a no/probably-no response to 4.8 if no
sensitivity analysis was attempted; that response must be accompanied
by the fact that no such analysis was done. It is not evidence that
an unperformed analysis established stability.
\end{hbmethod}

\begin{hbmethod}{def:p15-missing-robme-judgment}{Follow the suggested paths and justify the judgement}
Group yes/probably yes and no/probably no when following the tool's
suggested algorithm, but retain the original answers and reasons.
First determine the within-study situation:
\begin{enumerate}
\item Known selectively unavailable results with a likely notable
      impact (4.1 yes and 4.2 yes/probably yes) lead to high risk.
\item No missing or potentially missing results (4.1 and 4.3 both no)
      allow low risk if there is no potential for missing studies
      (4.5 no). If there is potential, ask 4.6. When eligible results
      are unlikely, proceed to the observed-pattern branch; when
      likely, proceed to the stronger across-study concern branch.
\item Missing or potentially missing results with no likely notable
      impact (4.2 and 4.4 no/probably no/not applicable) proceed
      through 4.5 and, where triggered, 4.6. No across-study potential,
      or unlikely eligibility of missing results, leads to the
      observed-pattern branch; likely eligibility leads to the
      stronger across-study concern branch.
\item An unclear notable impact of known missing results (4.2 no
      information), or a likely/unclear notable impact of potentially
      missing results (4.4 yes/probably yes/no information), retains
      within-study uncertainty. Even without strong additional
      across-study evidence, the suggested judgement is some concerns.
\end{enumerate}
In the observed-pattern branch for otherwise reassuring within-study
evidence, 4.7 no/probably no supports low risk. If 4.7 is
yes/probably yes, 4.8 yes/probably yes supports high risk and
4.8 no/probably no gives some concerns.
With retained within-study uncertainty but no likely eligible
unidentified results, either 4.7 or 4.8 yes/probably yes supports
high risk; two negative responses leave some concerns.
In the stronger across-study branch (4.5 yes and 4.6 yes/probably yes),
either 4.7 or 4.8 yes/probably yes supports high risk; otherwise
the suggested judgement is some concerns.

These are assessment paths, not a numerical sum of affirmative answers
or a proof of the actual bias. The source figure does not display an
exit for every permissible not-applicable answer at the pattern node.
Preserve that answer and explain the final judgement from the evidence;
do not silently convert inability to assess a pattern into a negative
finding. Record the tool's proposed judgement, any departure and its
basis, and optionally the anticipated direction: favouring either
intervention, toward or away from the null, or unpredictable.
A low ROB-ME judgement concerns missing evidence for this result; it
does not certify the studies' conduct, reporting-choice integrity,
transportability or absence of heterogeneity
\cite[fig.~2, table~1, pp.~13--15]{ROBME2023Cribsheet}
\cite[sec.~13.3.5]{PageHigginsSterne2024Cochrane13}.
\end{hbmethod}

\subsection{Sensitivity with an explicit missing-evidence assumption}

\begin{hbdefinition}{def:p15-missing-weight-bounds}{Known missing weight and an effect range}
Let the available estimates have total fixed positive weight \(W_o\)
and weighted mean \(M_o\). Identified missing results \(y_j\)
have specified positive weights \(w_j\) and ranges
\(L_j\leq y_j\leq U_j\). Put \(W_m=\sum_jw_j\).
These are explicitly assumed numerical values or ranges; sample size
alone does not determine a missing result's inverse-variance weight.
If the missing standard errors or effects are unknown, vary their
assumed values rather than labeling them observed.
\end{hbdefinition}

\begin{hbtheorem}{thm:p15-missing-weight-bounds}{proved=annex:p15.missing.bounds}{Missing-weight bounds and a sign tipping point}
Under these fixed-weight assumptions, the completed summary lies in
\[
 \left[\frac{W_oM_o+\sum_jw_jL_j}{W_o+W_m},
       \frac{W_oM_o+\sum_jw_jU_j}{W_o+W_m}\right],
\]
with both endpoints attainable. For \(W_m>0\), its shift from
\(M_o\) is \(W_m(M_m-M_o)/(W_o+W_m)\), where
\(M_m=\sum_jw_jy_j/W_m\).
It is at most a chosen threshold \(d\) exactly when
\[
                 \sum_jw_jy_j\leq d(W_o+W_m)-W_oM_o.
\]
For three available estimates each equal to two with unit weights,
and two missing unit-weight estimates each allowed in \([-4,0]\),
the completed mean ranges from \(-2/5\) to \(6/5\).
It is nonpositive exactly when the two missing estimates sum to
at most \(-6\). No confidence coverage is claimed for these
deterministic sensitivity bounds.
\end{hbtheorem}

Compare predeclared common-effect and random-effects calculations,
and justified restrictions to larger or more credible studies.
The changed estimand and applicability of a restricted study set
must be stated. Inverse-variance random-effects weighting gives
less relative dominance to precise studies, so a small-study effect
can become more influential; it is not a selection correction.

\begin{hbtheorem}{thm:p15-missing-small-weight}{proved=annex:p15.missing.small-weight}{The exact two-study weight shift}
For two fixed estimates \(y_L,y_S\), with variances
\(0<v_L<v_S\), put
\[
 M(u)=\frac{(v_S+u)y_L+(v_L+u)y_S}{v_L+v_S+2u},
                         \qquad u\geq0.
\]
Then
\[
 M'(u)=\frac{(v_S-v_L)(y_S-y_L)}{(v_L+v_S+2u)^2},\qquad
 \lim_{u\to\infty}M(u)=(y_L+y_S)/2.
\]
Thus the less precise result gains relative weight as the common
heterogeneity allowance increases. With \((y_L,y_S)=(0,3)\)
and \((v_L,v_S)=(1,9)\), the means are \(3/10\) at zero
and \(5/6\) at \(u=4\). Changing weights of exclusively
positive available estimates cannot produce a nonpositive weighted
mean; it cannot discover the missing negative results of the preceding
example.
\end{hbtheorem}

Selection-model sensitivity assumes an availability law depending on
effect, significance, precision or their joint structure. The selection
identity above shows what specified probabilities accomplish and what
unknown probabilities leave unidentified. Vary plausible severity,
missing-study count and omitted-result assumptions; do not call one
fitted correction the recovered truth. An assumed selection mechanism
can be misleading when real heterogeneity or within-study bias produced
the observed pattern.
Regression adjustments extrapolate an effect-size/variance relation
to greater precision, sometimes formally to zero variance. Such
extrapolation needs sufficient studies and a supported functional
form; when all studies are small, its intercept may lie beyond every
observed effect. A fitted line does not establish that its slope is
caused by missing results. Report the substantive range of assumptions
and resulting changes in magnitude, sign and uncertainty
\cite[sec.~13.3.4.6]{PageHigginsSterne2024Cochrane13}.

\section{Funnel asymmetry and alternative explanations}
\label{sec:p15-missing-funnel}

A funnel plot places study effect estimates horizontally and precision
or size vertically, commonly a standard error on a reversed vertical
scale so that more precise studies lie toward the top. Ratio effects
use a logarithmic horizontal scale. Sample size alone does not encode
the number of events, outcome variability or design dependence.
Under a common-effect normal model with known centre, reference flanks
\(\theta\pm1.96s\) describe individual approximate 95 percent
sampling ranges. They are not a simultaneous guarantee for all studies,
and fitted centres, heterogeneity or correlated estimate/SE construction
change their interpretation.

The selected source lists several causes of small-study effects:
result-dependent availability; greater bias in small-study conduct or
reported-result selection; true effect differences across populations
or implementation; mathematical association of an effect estimate with
its standard error; and chance. A complete inception cohort with fully
available eligible results can have an asymmetric funnel because
its smaller studies differ substantively from its larger studies.
That asymmetry alone does not raise a missing-evidence diagnosis
\cite[secs.~13.3.4.2--13.3.4.5, 13.3.5]{PageHigginsSterne2024Cochrane13}.

\begin{hbtheorem}{thm:p15-missing-funnel}{proved=annex:p15.missing.funnel}{One asymmetric plot compatible with three mechanisms}
Consider stipulated rounded estimates from ten studies. Five have
SE \(1/10\), with displayed effects
\(-1/10,-1/20,0,1/20,1/10\). Five have SE one, with displayed
effects \(8/5,9/5,2,11/5,12/5\).
For any sufficiently small positive rounding half-width, the event
of all estimates lying in their displayed intervals has positive
probability under each of these models:
\begin{enumerate}
\item All studies are available, all true effects are zero, and
      estimates have independent normal errors with the given SEs.
\item All studies are available, the precise studies have true effect
      zero and the less precise studies true effect two, with the
      same independent normal error laws.
\item True effects are zero, all precise results are available, and
      less precise studies' normal results are made available only
      when their estimate exceeds one.
\end{enumerate}
The first explains this realised pattern by chance, the second permits
true heterogeneity, and the third uses selective availability.
Their likelihoods need not agree. Compatibility of the observations
does not establish a unique cause of their asymmetry.
\end{hbtheorem}

Contour lines at chosen significance thresholds add information about
where apparently missing results would fall relative to a stated null.
A gap in a nonsignificant, unfavourable region can strengthen a
selection explanation; a gap among highly significant favourable
results can weaken that particular explanation. The contours do not
enumerate unseen studies or exclude the other mechanisms.

Funnel-asymmetry tests assess an effect-size/precision association,
not a publication mechanism directly. The selected chapter advises
at least ten studies as a rule of thumb, adequate variation in
standard errors, visual and contextual interpretation, and an
effect-scale-appropriate test. Low power means that a negative test
does not exclude selection. Ten studies do not guarantee good power.
The original Egger procedure can have artefactual problems for odds
ratios and standardized mean differences; the source distinguishes
other methods for these settings. A significance threshold from a
generic regression is not automatically valid for every effect scale
\cite[sec.~13.3.4.4]{PageHigginsSterne2024Cochrane13}.

\section{Chapter versions and the interpretation of combined evidence}
\label{sec:p15-missing-interpretation}

Keep the review question, eligible study denominator and result definition,
availability matrix, source-search record, ROB-ME version and answers,
judgement reasons, sensitivity assumptions and actual calculated results
together. If evidence is insufficient for a trustworthy pooled conclusion,
describe the available results and missing eligible evidence explicitly.
Numerical precision, absence of significant asymmetry and agreement
between chosen sensitivity analyses do not establish a complete corpus.
Nonselective unavailability can still limit the target and its applicability;
a low rating in one bias domain is not a general validity certificate.

The preceding synthesis treatment uses the distinct November~2024
Chapter~10 presentation for its pooling methods. This missing-evidence
treatment uses Chapter~13's August~2024 presentation and the October~2023
ROB-ME operating documents. Dates and versions identify different
methodological objects; a current webpage or the title of a whole
handbook does not justify combining incompatible variants.
The selected chapter's summary recommends considering the impact of
non-reporting in the review's conclusions; its MECIR expectation C73
classifies this investigation as highly desirable.
Thorough search and attempted recovery reduce opportunities for missing
evidence but do not, by themselves, prove their absence
\cite[sec.~13.4]{PageHigginsSterne2024Cochrane13}.
